% ===========================================================================
%   DROPPING A QUANTUM TEMPORARY -- supplementary material.
%
%   The paper's appendices, as a separate PDF, so that full.tex stands alone.
%   What goes here is VERIFICATION, never exposition: a paper that needs its
%   appendix to be believed has not been written yet.
%
%   A \cref from full.tex into this document DOES NOT RESOLVE.  Say "the
%   supplement, Appendix A" in prose there.
%
%   Statements shared with the body are macros in statements.tex.  Appendix A
%   is ordered as the body is; Appendix B is the catalogue.
% ===========================================================================

\documentclass[acmsmall,screen,review,anonymous,nonacm]{acmart}

\usepackage{amsthm}
\usepackage{enumitem}
\usepackage{array}
\usepackage{framed}

\usepackage{mathpartir}

\usepackage{cleveref}

\setlength{\emergencystretch}{2em}
\hyphenpenalty=1000
\sloppy

\theoremstyle{acmplain}
\newtheorem{theorem}{Theorem}[section]
\newtheorem{proposition}[theorem]{Proposition}
\newtheorem{lemma}[theorem]{Lemma}
\newtheorem{corollary}[theorem]{Corollary}
\newtheorem{fact}[theorem]{Fact}

\theoremstyle{acmdefinition}
\newtheorem{definition}[theorem]{Definition}
\newtheorem{example}[theorem]{Example}
\newtheorem{remark}[theorem]{Remark}

\makeatletter
\@ifpackageloaded{cleveref}{%
  \newcommand{\crefsharedcounter}[1]{%
    \@ifundefined{#1}{}{%
      \expandafter\let\csname clv@#1\expandafter\endcsname\csname #1\endcsname
      \expandafter\let\csname endclv@#1\expandafter\endcsname\csname end#1\endcsname
      \expandafter\def\csname #1\endcsname{\crefalias{theorem}{#1}\csname clv@#1\endcsname}%
      \expandafter\def\csname end#1\endcsname{\csname endclv@#1\endcsname}%
    }%
  }%
  \crefsharedcounter{proposition}%
  \crefsharedcounter{lemma}%
  \crefsharedcounter{corollary}%
  \crefsharedcounter{fact}%
  \crefsharedcounter{definition}%
  \crefsharedcounter{example}%
  \crefsharedcounter{remark}%
}{}
\makeatother

\crefname{fact}{fact}{facts}
\Crefname{fact}{Fact}{Facts}

\renewcommand{\qedsymbol}{\ensuremath{\square}}

\definecolor{accent}{HTML}{1F4E79}
\definecolor{accentwarm}{HTML}{9A5B00}

\newcommand{\rulename}[1]{{\normalfont\scshape\footnotesize\color{accent}#1}}

\mprset{andskip=1.6em plus 0.4fil minus 0.4em, sep=0.9em}
\newcommand{\ruleskip}{\vskip 0.25\baselineskip}

\ifPDFTeX\else
  \setmonofont{Iosevka Custom}[Scale=MatchLowercase,RawFeature={-liga,-calt,-clig,-dlig}]
\fi

\newlist{steps}{enumerate}{1}
\setlist[steps]{label=$\langle$\arabic*$\rangle$,ref=$\langle$\arabic*$\rangle$,
  leftmargin=2.1em,labelsep=0.5em,align=left,itemsep=0.5pt,topsep=1pt,parsep=0pt}

\newcommand{\code}[1]{\texttt{\small #1}}

\setlength{\heavyrulewidth}{0.9pt}
\setlength{\lightrulewidth}{0.4pt}
\setlength{\aboverulesep}{3pt}
\setlength{\belowrulesep}{3pt}

\setlength{\parskip}{0pt}

\setlist[itemize]{leftmargin=1.3em, itemsep=0.25\baselineskip, topsep=0.4\baselineskip, parsep=0pt, partopsep=0pt}
\setlist[enumerate]{leftmargin=1.6em, itemsep=0.25\baselineskip, topsep=0.4\baselineskip, parsep=0pt, partopsep=0pt}
\setlist[description]{leftmargin=0pt, itemsep=0.3\baselineskip, topsep=0.4\baselineskip, partopsep=0pt, font=\normalfont\itshape}

\newcommand{\C}{\mathbb{C}}
\newcommand{\Q}{\mathbb{Q}}
\newcommand{\Z}{\mathbb{Z}}
\newcommand{\Hilb}{\mathcal{H}}
\newcommand{\ket}[1]{\lvert #1 \rangle}
\newcommand{\bra}[1]{\langle #1 \rvert}
\newcommand{\proj}[1]{\lvert #1 \rangle\!\langle #1 \rvert}
\newcommand{\id}{\mathrm{id}}

\newcommand{\den}[1]{[\![#1]\!]}

\newcommand{\emptyctx}{\cdot}

\newcommand{\drop}{\texttt{drop}}
\newcommand{\lift}{\texttt{lift}}
\newcommand{\new}{\texttt{new}}
\newcommand{\sep}{\mathrm{sep}}

\newcommand{\jty}[3]{#1 \vdash #2 : #3}

\newcommand{\jrule}[3]{\inferrule*[right={\rulename{#1}}]{#2}{#3}}
\newcommand{\jaxiom}[2]{\inferrule*[right={\rulename{#1}}]{\ }{#2}}

\newenvironment{ruledisplay}
  {\par\ruleskip\begin{mathpar}}
  {\end{mathpar}\par\ruleskip}

\definecolor{planink}{HTML}{9A5B00}
\newcounter{plancount}

\newenvironment{planned}[1]
  {\par\vskip 0.45\baselineskip%
   \stepcounter{plancount}%
   \def\FrameCommand{{\color{planink}\vrule width 1.6pt}\hspace{0.6em}}%
   \MakeFramed{\advance\hsize-\width\FrameRestore}%
   \noindent{\color{planink}\sffamily\bfseries\scriptsize Planned\hspace{0.9em}}%
   \small\ignorespaces}
  {\endMakeFramed\par\vskip 0.45\baselineskip\noindent\ignorespaces}

\newcommand{\exRunning}{%
\begin{example}[the running examples]\label{ex:three}
\leavevmode
\begin{center}\small
\begin{tabular}{@{}lll@{}}
\toprule
(a) \emph{superpose, then drop} & (b) \emph{superpose, then read} & (c) \emph{copy, measure, one joint gate} \\
\midrule
\texttt{prog a (v : [2]) : [2] \{} & \texttt{prog b (v : [2]) : [2] \{} & \texttt{prog c (v : [2]) : [2] \{} \\
\texttt{\ \ new t : [2] @ |0>;}    & \texttt{\ \ new t : [2] @ |0>;}    & \texttt{\ \ new t : [2] @ |0>;} \\
\texttt{\ \ H t;}                  & \texttt{\ \ H t;}                  & \texttt{\ \ H v;\ \ CX v t;} \\
                                   & \texttt{\ \ CX t v;}               & \texttt{\ \ p = lift[\{0\}] t;} \\
                                   &                                    & \texttt{\ \ V t v;} \\
\texttt{\ \ drop t;}               & \texttt{\ \ drop t;}               & \texttt{\ \ drop t;} \\
\texttt{\ \ return v \}}           & \texttt{\ \ return v \}}           & \texttt{\ \ return v \}} \\
\bottomrule
\end{tabular}
\end{center}
\end{example}}

\newcommand{\stmtChannel}{%
\begin{proposition}[the summed channel does not decide]\label{prop:channel}
There are two well-typed programs on a temporary qubit and a visible qubit,
each with two reachable records, that have the same summed channel $\Phi_0 +
\Phi_1$ and the same weight on every branch at every input, and opposite
verdicts on every branch.
\end{proposition}}

\newcommand{\stmtProductwise}{%
\begin{proposition}[a product at every input]\label{prop:productwise}
Let $K : \den{\Omega} \to \Hv \otimes \Hrest$ be linear, with $n = \dim\den{\Omega}$,
and suppose $Kx$ is a product vector, or zero, for every $x$. Then $K$ has one of
two shapes, and both exactly when $K$ has rank at most one.
Here the language spaces $\Hv$ and $\Hrest$ have positive dimension;
the empty register tuple denotes $\C$.
\begin{description}
\item[the temporary is fixed] $K = g \otimes W$ for a vector $g \in \Hv$
  and a linear $W : \den{\Omega} \to \Hrest$ --- the rest varies, the temporary
  does not;
\item[the rest is fixed] $Kx = Gx \otimes w$ for a fixed nonzero
  $w \in \Hrest$ and a linear $G : \den{\Omega} \to \Hv$ --- the temporary
  varies, the rest does not.
\end{description}
The first shape is what \cref{def:site} asks of a site map. A map of
the second shape is $g \otimes W$ exactly when $G$ has rank at most one, so
such an isometry with $n \ge 2$ never is.
\end{proposition}}

\newcommand{\stmtGram}{%
For a site map $K:\den\Omega\to\Hv\otimes\Hrest$, including a map on a measurement record, put
\[
 d=\dim\Hv\ge1,\qquad n=\dim\den\Omega\ge1,\qquad
 K_i=(\bra i\otimes I)K,\qquad
 \Gamma_{ij}=\operatorname{tr}(K_i^{\dagger}K_j).
\]
The Gram matrix $\Gamma$ records the pairwise inner products of the
operators $K_i$. No isometry hypothesis is needed in the proposition.
\begin{proposition}[the Gram test]\label{prop:gram}
\leavevmode
\begin{enumerate}[beginpenalty=10000]
\item \emph{Fixed block.} At a site without loans this is the safety test:
the following are equivalent:
\[
 K=g\otimes W\text{ for a unit }g
 \quad\Longleftrightarrow\quad \operatorname{rank}\Gamma\le1
 \quad\Longleftrightarrow\quad
 \Gamma^2=(\operatorname{tr}\Gamma)\Gamma.
\]
\item \emph{Least block dimension.} The map
$\rho\mapsto\sum_iK_i\rho K_i^{\dagger}$ can be represented using a
block of dimension $r\ge1$ exactly when
$r\ge\max(1,\operatorname{rank}\Gamma)$.
\item \emph{Changing gates.} A unitary on the rest leaves $\Gamma$
unchanged. A unitary $U$ on the block changes it to
$\bar U\Gamma U^{\top}$.
\end{enumerate}
\end{proposition}}

\newcommand{\stmtScalar}{%
\begin{proposition}[the block may learn nothing and the drop be unsafe]\label{prop:scalar}
Let $K=\sigma_s$ be the isometric map of a measurement-free site with
no loans in its block. Use the operators $K_i$ of \cref{prop:gram}.
The block has the same reduced state at every unit input exactly when
\[
 K_i^{\dagger}K_j=c_{ij}I\qquad\text{for all }i,j.
\]
Under this condition the drop is safe exactly when $c$ has rank one.
Then $c_{ij}=\bar g_i g_j$ for a unit certificate $g$. In contrast,
the program below has the constant mixed state $c=I/2$ and is unsafe.
\end{proposition}}

\newcommand{\stmtCliffordT}{%
\begin{theorem}[the Clifford$+T$ coefficient ring splits]\label{thm:cliffordT}
Let $\mathcal R = \Z[1/\sqrt2, i]$ and $d \ge 1$.
\begin{enumerate}
\item Every rank-one projection on $\C^{d}$ with entries in $\mathcal R$ is
  $zz^{\dagger}$ for a vector $z$ over $\mathcal R$.
\item For unit $g$, the ray $[g]$ carries a unit vector over $\mathcal R$ exactly
  when its projection $gg^{\dagger}$ has entries in $\mathcal R$.

  At a safe temporary site over $\mathcal R$, suppose the surviving weight
  at some basis input is a unit of the ring. The certificate then carries
  a unit vector over $\mathcal R$.

  A measurement-free temporary site over $\mathcal R$ has a witness over
  $\mathcal R$ exactly when it is safe: its record has weight $1$.
\end{enumerate}
\end{theorem}}

\newcommand{\stmtRings}{%
\begin{proposition}[the restricted rings]\label{prop:rings}
Over $\Z[1/2, i]$ and over $\Z[i/\sqrt2]$, every rank-one projection with
entries in the ring, in every dimension, is $zz^{\dagger}$ for a $z$ over the
ring. Over $\Z[1/2]$ the projection onto $\ket+$ is not, and over
$\Z[1/\sqrt2]$ the projection onto $\cos\tfrac{\pi}{8}\ket0 +
\sin\tfrac{\pi}{8}\ket1$ is not.
\end{proposition}}

\newcommand{\seqop}{\mathbin{+\mkern-8mu+}}
\newcommand{\Hhid}{\Hilb_{\mathsf{hid}}}

\newcommand{\stmtTypeSafety}{%
\begin{theorem}[type safety]\label{thm:typesafety}
Let $\vdash \langle \Delta, \sigma \parallel M\rangle : A$ be a well-typed
configuration of the measurement-free machine.
\begin{enumerate}
\item \emph{Progress.} Either $M$ is \texttt{return }$v$, and the
  configuration is final, or exactly one rule applies to it, and the successor is
  determined once the name \textsc{Step-New} allocates is fixed.
\item \emph{Preservation.} The step yields a well-typed configuration at
  the same type $A$.
\item \emph{Termination.} Every run from the configuration reaches a final
  configuration in exactly $|M|$ steps, where $|M|$ counts its statements
  and does not count \texttt{return}.
\item \emph{Adequacy.} The run from $\langle \Omega, x \parallel M\rangle$,
  the input registers allocated and the names it allocates fresh for $M$, ends in the
  state $\den{M}\,x$.
\end{enumerate}
\end{theorem}}

\newcommand{\stmtDesugar}{%
\begin{lemma}[desugaring]\label{lem:desugar}
Suppose $\Omega_1\vdash M_1:A$,
$\Omega_{\bar y}\vdash\bar y:A$, and
$\Omega_2,\Omega_{\bar y}\vdash M_2:B$.
The three contexts are pairwise disjoint. Every name bound in $M_1$
(register or parameter) is fresh for $M_2$, $\bar y$ and $\Omega_2$.
Every name bound in $M_2$ is fresh for $\Omega_1$ and $\bar y$.
\begin{enumerate}
\item $\Omega_1, \Omega_2 \vdash M_1 \seqop_{\bar y} M_2 : B$.
\item Every run of $M_1$ to \texttt{return }$v$, followed by the run of
  $M_2[v/\bar y]$ from the state it reaches, is a run of $M_1 \seqop_{\bar y} M_2$:
  same record and state. In the measurement-free case,
  \[
  \den{M_1\seqop_{\bar y}M_2}
  =(\id_{\Hhid^{M_1}}\otimes\den{M_2})
    \circ(\den{M_1}\otimes\id_{\den{\Omega_2}}).
  \]
\item $|M_1 \seqop_{\bar y} M_2| = |M_1| + |M_2|$.
Here $|M|$ counts statements along the longest branch of $M$:
\texttt{return} counts zero, each other statement counts one, and an
\texttt{if} counts one plus the larger count of its two arms.
\end{enumerate}
The disjointness of $\Omega_{\bar y}$ from $\Omega_1$ is not decoration:
without it the componentwise renaming can clone a register. The converse of (1) is
false --- the continuation may name a register of $\Omega_1$ directly.
\end{lemma}}

\newcommand{\stmtCompile}{%
\begin{theorem}[compilation]\label{thm:compile}
Let $M$ be a measurement-free program in the compilation fragment.
Assume each safe site $s_1,\ldots,s_r$ is emittable. Write $E_M$ for the
emitted circuit's map and $R^{s_j}$ for the ideal reclaiming gate at $s_j$.
\begin{enumerate}
\item \emph{Full state.} There is one scalar $z$, $|z|=1$, such that
\[
 E_M=z\,(R^{s_1}\otimes\cdots\otimes R^{s_r}\otimes I)\circ\den M.
\]
The output factors are grouped by dropped block, in site order, followed
by the remaining factors. The scalar $z$ is independent of every input,
including the states of the loans; it is the product of the phases
omitted during emission.
\item \emph{Returned output.} The emitted circuit and $M$ have the same
reduced state on the returned registers at every input. If no reclaiming
gate is inserted, their full maps differ only by $z$.
\item \emph{Reusable registers.} Each reclaimed temporary is in its
allocation state. Each reclaimed loan is in its original state, with
its correlations with the caller preserved. These are equalities of
states, so the global phase in (1) does not affect them.
\end{enumerate}
\end{theorem}}

\newcommand{\Hv}{\Hilb_{v}}
\newcommand{\Hrest}{\Hilb_{\mathsf{rest}}}
\newcommand{\Hloan}{\Hilb_{\mathrm{loan}}}
\newcommand{\Hother}{\Hilb_{\mathsf{other}}}
\newcommand{\pfx}[2]{#1|_{#2}}

\newcommand{\stmtSite}{%
\begin{definition}[the site, and safety at it]\label{def:site}
Let $\jty{\Omega}{M}{A}$ be a program.
\begin{enumerate}[beginpenalty=10000]
\item A \emph{site} $s$ of $M$ is an occurrence of $\drop\,v$ in it. Its
  \emph{block} is the ordered list of registers $v$ names --- none for $v = ()$, a site
  with $\Hv = \C$, safe with the unit scalar as certificate;
  $\Hv$ is the tensor product of
  their spaces in that order. The space $\Hrest$ is the tensor product of all remaining
  register spaces: inputs in context order, followed by allocations,
  including previously dropped registers (\cref{sec:lang:dynamics}).
\item The \emph{site map} $\sigma_s : \den{\Omega} \to \Hv \otimes \Hrest$
  sends an input vector $\xi\in\den\Omega$ to the joint state of all
  registers just before the drop. The block factors are listed first;
  no physical gate is applied by this regrouping.
\item First suppose none of the registers named by this drop is marked
  \texttt{loan}. They may be allocated temporaries or unmarked inputs.
  The site is \emph{safe} when its site map
  factors,
  \[ \sigma_s = g \otimes W, \]
  for a unit vector $g \in \Hv$, the site's \emph{certificate}, and a linear
  $W : \den{\Omega} \to \Hrest$: the block comes out in the one state $g$
  whatever the input.
\item If the drop also names registers marked \texttt{loan}, they must be
  returned in their unknown original states. Group those input factors as $\den{\Omega} =
  \Hloan \otimes \den{\Omega'}$, the block's loans and the rest, and the site
  is \emph{safe} when
  \[ \sigma_s = B \otimes W \]
  for an isometry $B : \Hloan \to \Hv$, the site's \emph{block factor}, and a
  linear $W : \den{\Omega'} \to \Hrest$. Clause (3) is the case $\Hloan = \C$,
  $B = g$.
\end{enumerate}
\end{definition}}

\newcommand{\stmtModular}{%
\begin{proposition}[safety is modular]\label{prop:modular}
Let $\sigma_s$ be the linear map of a site as in \cref{def:site}.
Let $\Hilb_c\ne0$ be a client register on which the program does nothing.
Then
\[
 \sigma_s\otimes I_c=B\otimes W'
 \text{ for some }W'
 \quad\Longleftrightarrow\quad
 \sigma_s=B\otimes W\text{ for some }W.
\]
The same isometric block factor $B$ works on both sides, including when
the client is entangled with the inputs. For a block without loans,
read $B=g$.
\end{proposition}}

\newcommand{\stmtIsoExt}{%
\begin{lemma}[isometry extension]\label{lem:isoext}
Let $\Hilb$ be a finite-dimensional complex Hilbert space. Every isometry
$V : S \to \Hilb$ out of a subspace $S \subseteq \Hilb$ is the restriction of
a unitary of $\Hilb$.
\end{lemma}}

\newcommand{\stmtCriterion}{%
\begin{theorem}[the criterion, and contextual equivalence]\label{thm:criterion}
In the Setting, the following are equivalent.
\begin{enumerate}
\item $s$ is safe.
\item $\sigma_s = g \otimes W$ for a unit $g \in \Hv$ and an isometry
  $W : \den{\Omega} \to \Hrest$.
\item Some weakening witness $\sep(R,W)$ has denotation $\sigma_s$.
\item No closed context tells $\pfx{M}{s}$ from some $\sep^{\circ}(R,W)$.
\end{enumerate}
Two prefixes no closed context tells apart have the same certificates.
\end{theorem}}

\newcommand{\stmtComposite}{%
\begin{proposition}[the composite temporary]\label{prop:composite}
Let $M$ have temporaries $t_1,\ldots,t_n$, each dropped once on the
chosen record, and let $\Phi$ be its final state map on that record.
For clause (1), assume that the tail after the drop of $t_k$ is
injective on the span of the row vectors
$(\bra{e_i}\otimes I)\sigma_s(x)$, as $i$ and $x$ vary. This holds if
that tail contains no lifts. For a lift, the condition says that its
projection kills no nonzero vector in the relevant span.
Call $t_k$ \emph{separately
constant} when $\Phi(x)$ lies in $g_k \otimes (\text{the other factors})$ for
one unit $g_k \in \Hilb_{t_k}$ and every input $x$.
\begin{enumerate}
\item The site of $\drop\,t_k$ is safe if and only if $t_k$ is separately
  constant.
\item Every $t_k$ is separately constant if and only if the composite site
  $\drop\,(t_1, \ldots, t_n)$, placed at the end after deleting the
  individual drops, is safe with
  a certificate carrying the product $g_1 \otimes \cdots \otimes g_n$.
\end{enumerate}
\end{proposition}}

\newcommand{\stmtPhase}{%
\begin{lemma}[equal answers, proportional maps]\label{lem:phase}
Let two prefixes with every $\drop$ deleted have the same context and type,
so that their denotations $\sigma, \sigma' : \den{\Omega} \to \Hilb$ are
linear maps into one space. If no closed context tells them apart, then
$\sigma' = c\,\sigma$ for a scalar $c$, of modulus one once $\sigma \neq 0$.
Conversely, no closed context tells apart two prefixes with
$\sigma' = c\,\sigma$ and $|c| = 1$.
\end{lemma}}

\newcommand{\stmtReuse}{%
\begin{proposition}[reclaiming temporaries]\label{prop:reuse}
In the Setting above, define
\[
 R^{\psi}=\begin{cases}
 I,&w=0,\\
 I-2ww^{\dagger}/\lVert w\rVert^2,&w\ne0.
 \end{cases}
\]
\begin{enumerate}
\item \emph{Exchange.} $R^{\psi}$ is the unique linear map that exchanges
$\tilde g$ and $\psi$ and fixes their common orthogonal complement.
It is a unitary involution.
\item \emph{No gate needed.} $R^{\psi}=I$ exactly when $\tilde g=\psi$,
equivalently when $g$ lies on the allocation ray.
\item \emph{Restoration.} At every input, applying $R^{\psi}$ to the
block puts it in $\psi$ and leaves the reduced state of the rest unchanged.
\end{enumerate}
\end{proposition}
A dropped block may contain both loans and temporaries. A loan must be
returned in the state the caller supplied; a temporary must be returned
to its declared allocation state.

At a safe site, the block factor
$B:\Hloan\to\Hv$ of \cref{def:site}(4) describes where the program has
put the loans' information. For a loan input $\ell$, the block contains
$B\ell$. Let $J$ describe the desired result: $J\ell$ places the unchanged
loans beside the temporaries in their allocation state $\psi$.
If loans precede temporaries in the tuple, $J\ell=\ell\otimes\psi$;
otherwise these factors are placed in the tuple's order.

\begin{proposition}[returning loans]\label{prop:loanreturn}
For the safe block above, there is a unitary $T$ on the block such that
\[
                    T(B\ell)=J\ell\qquad\text{for every loan input }\ell.
\]
Thus the same gate restores every possible loan state and every
temporary's allocation state. It also preserves the loans' correlations
with registers kept by the caller.

Only the action of $T$ on the image of $B$ is prescribed. Different
unitaries may extend that action to the rest of the block space.
If the block contains only loans, $B$ is square and unitary:
$J=I$ and $T=B^{\dagger}$. In this case $T$ is the identity up to phase
exactly when $B$ is.
\end{proposition}}

\newcommand{\stmtReuseSound}{%
\begin{proposition}[reuse is sound]\label{prop:reusesound}
Extend the machine of \cref{fig:machine} with a pool: at a safe site,
\textsc{Step-Drop-Reclaim} applies $R^{s}$ to the block and moves its
registers to the pool instead of hiding them (loans are returned by the gate of \cref{prop:loanreturn}
and are never pooled);
\textsc{Step-New-Reuse} serves
a \texttt{new } $t : [d] \mathbin{@} \psi$ from a pooled register standing in $\psi$
instead of allocating a fresh one.
\begin{enumerate}
\item Every run of the pooled machine is simulated by a run of the machine of
  \cref{fig:machine}: the same records, the same weight at every input, and
  the same reduced state on the returned registers, on every record.
\item The pooled machine allocates at most as many registers, and strictly
  fewer whenever a reclaimed register is served.
\item If reclamation is allowed at an unsafe site, reuse can be unsound. Extend program (b) by
  allocating a fresh qubit and returning it with $v$. If that allocation
  instead reuses the dropped block, no gate on the block alone makes the
  returned joint state agree at every input.
\end{enumerate}
\end{proposition}}

\newcommand{\stmtTypeSafetyLift}{%
\begin{theorem}[type safety, with the two forms]\label{thm:typesafetylift}
\Cref{thm:typesafety} holds of the five-rule machine, with two clauses
read again and two unchanged.
\begin{enumerate}[beginpenalty=10000]
\item \emph{Progress} now reads: a well-typed configuration is final, or
  exactly one rule applies to it, and at \textsc{Step-Lift} both of that rule's
  labelled instances do --- a lift being the only labelled step.
\item \emph{Termination} now reads along the branch taken: the run whose
  record is $o$ reaches a final configuration in at most as many steps as that branch
  has statements, an \texttt{if} counting as one, and so in at most $|M|$.
\item \emph{Preservation} stands word for word, and so does the
  hidden-register invariant of \cref{sec:lang:dynamics}: a lift binds its parameter
  and leaves its register live under its own name, so no register enters or leaves
  the context there.
\item \emph{Adequacy} stands with the run read along its record: the run
  from $\langle \Omega, x \parallel M\rangle$ whose record is $o$ ends in the state
  $E_o(M)\,x$.
\end{enumerate}
\end{theorem}}

\newcommand{\stmtBlockLift}{%
\begin{remark}[a block lift is a derived form]\label{rem:blocklift}
Let $S\subseteq[d]$ be nonempty and proper. Define the indicator gate
\[
 U_S=\sum_{j\in S}\proj j\otimes X+
     \sum_{j\notin S}\proj j\otimes I.
\]
It leaves $x$'s level unchanged and flips a fresh qubit $a$ exactly when
that level belongs to $S$. In a signature containing $U_S$, replace
$p=\lift_S\,x\texttt{;}\ M$ by these steps:
\begin{enumerate}
\item Allocate $a$ in $\ket0$ and apply $U_S$ to $(x,a)$.
\item Set $p=\lift_{\{1\}}\,a$. Thus $p=1$ has the same meaning as in
      the original test: the level of $x$ belongs to $S$.
\item Drop $a$ and continue with $M$.
\end{enumerate}
To check the replacement, write the incoming state as
$\sigma=\sum_j\ket j_x\otimes\sigma_j$ and put
$\tau_1=\sum_{j\in S}\ket j\otimes\sigma_j$ and
$\tau_0=\sum_{j\notin S}\ket j\otimes\sigma_j$.
Step (1) gives $\tau_1\otimes\ket1_a+\tau_0\otimes\ket0_a$.
On outcome $b$, step (2) keeps
$\tau_b\otimes\ket b_a=((P_S^b)_x\sigma)\otimes\ket b_a$.
This is the original branch state with one extra, fixed factor.
The replacement therefore has the same type and per-record map after
omitting that factor. Its drop of $a$ is safe, and $X^b$ restores $a$
to $\ket0$ (\cref{prop:reuse}).
\end{remark}}

\newcommand{\stmtReuseRecord}{%
\begin{proposition}[the record chooses the reclaiming gate]\label{prop:reuserecord}
Let $s$ fire on record $o$ with $\sigma_s^o\ne0$, safe there with certificate
$[g_o]$ and block allocated in $\psi$, and let $R_o$ be the reclaiming gate
of \cref{prop:reuse} for that certificate.
\begin{enumerate}
\item $g_o$ and $R_o$ depend on $o$ only through the bits $o$ takes at the
  lifts that fire before $s$: two records agreeing on those have the same
  certificate and the same reclaiming gate.
\item One gate cannot serve all records: some well-typed $M$ has a
  site whose reclaiming gate is $I$ on one record and $X$ on the other,
  whether or not an \texttt{if} of $M$ reads the bit that decides.
\end{enumerate}
\end{proposition}}

\newcommand{\stmtCompileLeaf}{%
\begin{theorem}[compilation, one record at a time]\label{thm:compileleaf}
Let $M$ be in the compilation fragment, with every lifted register of
power-of-two dimension. Fix a record $o$. List the safe sites that fire
on this record as $s_1,\ldots,s_r$, in execution order. Assume that each
site's certificate or loan block factor is emittable, and write
$R_o^{s_j}$ for its ideal reclaiming gate.

Let $C_o$ be the corresponding branch map of the emitted circuit.
When comparing maps, omit any additional indicator qubits that the
emitter returns to $\ket0$. Order the remaining hidden factors by site
and drop-tuple order, as in \cref{thm:compile}. Then:
\begin{enumerate}
\item The compiled branch is the source branch followed by the chosen
  reclaiming gates:
  \[
  C_o=z_o\bigl(R_o^{s_1}\otimes\cdots\otimes R_o^{s_r}\otimes\id\bigr)E_o(M),
  \qquad |z_o|=1.
  \]
  The phase $z_o$ is independent of every input, including loan states.
\item The branch weight and the reduced output state agree with those
  of $E_o(M)$ for every input. If no gate is emitted, the entire
  branch maps agree up to $z_o$.
\item Each reclaimed temporary is in its declared allocation state,
  and each reclaimed loan is in its original state. The phase in
  clause (1) changes neither of these states.
\end{enumerate}
\end{theorem}}

\newcommand{\stmtRecordWitness}{%
\begin{theorem}[the criterion, one record at a time]\label{thm:perrecord}
Let $s$ be a site of $M$ whose block is a tuple of temporaries allocated in
the unit state $\psi$, and $o$ a record on which $s$ fires. The following are equivalent.
\begin{enumerate}
\item $s$ is safe on $o$.
\item $\sigma^{o}_s = g \otimes W_o$ for a unit $g \in \Hv$ and a contraction
  $W_o : \den{\Omega} \to \Hrest$.
\item A witness has map $\sigma_s^o$ on its \emph{kept record},
where its fresh qubit $a$ reads $0$, after omitting that fixed
$\ket0_a$ factor. The witness performs three steps:
prepare the block in $\psi$ and apply one gate $R$ to it;
allocate the other temporaries and $a$ in $\ket{0}$ and apply one gate $U$
to everything except the block; then lift $a$ against $\{0\}$.
\item No closed context, and already no lift-free one, tells $\pfx{M}{s}$ on
  $o$ from some such program on its kept record.
\end{enumerate}
\end{theorem}}

\newcommand{\stmtObsInvRecord}{%
\begin{corollary}[safety is an observational invariant, one record at a time]
\label{cor:obsinvrecord}
Fix a record $o$ and let two prefixes with every $\drop$ deleted have the
same context and type, and let no lift-free closed context tell them apart on $o$. Then
${\sigma'}^{o} = c\,\sigma^{o}$ for a scalar $c$, of modulus one once
$\sigma^{o} \ne 0$; hence their sites are both safe or both unsafe on $o$, and every
certificate of one is a certificate of the other.
\end{corollary}}

\newcommand{\stmtNoTextQubit}{%
\begin{theorem}[the verdict is record-dependent]\label{thm:notext}
Let $c$ be program (c) of \cref{sec:meas:programs}.
\begin{enumerate}
\item $c$ is well typed, and its two records are both reachable, each of
  weight $\tfrac12$ at every basis input.
\item Under the record taking $1$ at the lift, the temporary landed
  in $\{0\}$, and the site map is
  $\sigma_{\mathrm{in}}(x) = \tfrac{x_0 + x_1}{\sqrt2}\,\ket{0} \otimes \ket{0}$,
  and the drop is safe, with certificate $\ket{0}$.
\item Under the record taking $0$ the site map is
  $\sigma_{\mathrm{out}}(x) = \tfrac{x_0 - x_1}{2}\,(\ket{0} \otimes \ket{1}
   + \ket{1} \otimes \ket{0})$,
  and the drop is unsafe.
\end{enumerate}
\end{theorem}}

\newcommand{\stmtRelativisedSite}{%
\begin{theorem}[the criterion over a gate ring]\label{thm:relativised}
Let $\mathcal A$ be a gate ring, $M$ a measurement-free program over
$\mathcal A$, and $s$ a site of $M$ whose block is a tuple of temporaries
allocated in $\psi \in \mathcal A^{d}$.
\begin{enumerate}
\item If some weakening witness $\sep(R,W)$ with $R$ over $\mathcal A$ has
  denotation $\sigma_s$, then $s$ is safe with a certificate carrying the unit
  vector $R\psi$ over $\mathcal A$, whatever $W$ is.
\item If $\mathcal A$ has a column lemma (\cref{fact:column}) and $s$ is safe
  with a certificate carrying a unit vector over $\mathcal A$, then some
  weakening witness \emph{over $\mathcal A$} has denotation $\sigma_s$.
\end{enumerate}
\end{theorem}}

\newcommand{\stmtClassicalCostSite}{%
\begin{theorem}[the classical price]\label{thm:classicalcost}
Deciding whether a given site whose block contains no loans is safe on a
given record is coNP-complete, for programs whose gates are drawn from $\{X, \mathrm{CNOT},
\mathrm{Toffoli}\}$ on qubit registers and whose allocations are clean.
Hardness holds already for measurement-free programs with one site.
\end{theorem}}


\settopmatter{printacmref=false, printfolios=true}
\setcopyright{none}
\acmConference[Under review]{Submitted for double-blind review}{}{}
\acmYear{2026}

% Appendix A, B -- so a referee can tell at a glance which document a
% number came from.  Floats follow the sections for the same reason: a bare
% "Figure 1" here is indistinguishable from the body's Figure 1, and sits
% beside theorems already numbered A.n.
\renewcommand{\thesection}{\Alph{section}}
\counterwithin{figure}{section}
\counterwithin{table}{section}

\begin{document}

\title{Dropping a Quantum Temporary}
\subtitle{Supplementary material}

\author{Author Name}
\affiliation{%
  \institution{Indiana University}
  \city{Bloomington}
  \state{Indiana}
  \country{USA}}
\email{sdkadu@iu.edu}

\maketitle

\noindent\textit{The proofs the paper defers and the catalogue it summarises.
This document is not self-contained and is not meant to be read alone.}


% ===========================================================================
% A.  DEFERRED PROOFS
% ---------------------------------------------------------------------------
% Restate, then prove. The statement a reader sees here must be the SAME
% statement, so each subsection CALLS the macro from statements.tex and never
% retypes it. Theorem numbers restart in this document; the prose names each
% by its title. A boxed "Planned" says which proofs are not here yet.
% ===========================================================================
\section{Deferred proofs}\label{app:proofs}

The proofs in this appendix are the ones the body states in one sentence
each; nothing in them is used before it is stated. The conventions are the
body's. A map $K$ from the visible register to the joint register has
\emph{Kraus operators} $K_i = (\langle i| \otimes I)K$ and is \emph{drop-safe}
when $K = g \otimes W$, equivalently when the $K_i$ are pairwise
proportional; the zero map is drop-safe. A record is \emph{reachable} for a
command when the selected denotation along it is nonzero. The three programs of the body
are recalled first, since the proofs are checked on them.

\exRunning

\subsection{The lifted calculus in full, and its metatheory}\label{app:meta}

The body displays the grammar and three reductions of the command calculus.
Here are the rest. A command is typed in a \emph{parameter count} $n$, the
number of parameters bound so far, and a \emph{flag} $\ell \in
\{\mathsf{live}, \mathsf{dropped}\}$ for the temporary; the parameter bound by
the $n$-th lift is written $p_n$, and a boolean expression $e$ is built from
parameters, negation and conjunction. The rules are
\begin{ruledisplay}
  \jaxiom{Halt}{n;\ell \vdash \texttt{halt}}
\and
  \jrule{Gate}{U^{\top}U = I \\ n;\mathsf{live} \vdash c}
              {n;\mathsf{live} \vdash \texttt{gate}\ U\,;c}
\and
  \jrule{VGate}{V^{\top}V = I \\ n;\ell \vdash c}
               {n;\ell \vdash \texttt{vgate}\ V\,;c}
\and
  \jrule{Drop}{n;\mathsf{dropped} \vdash c}
              {n;\mathsf{live} \vdash \drop\ t\,;c}
\and
  \jrule{Lift}{S \subseteq [d] \\ n+1;\mathsf{live} \vdash c}
              {n;\mathsf{live} \vdash
               \texttt{let } p_n = \lift(\meas\ t \in S)\texttt{ in } c}
\and
  \jrule{If}{\mathrm{params}(e) \subseteq \{p_0, \ldots, p_{n-1}\} \\
             n;\ell \vdash c_1 \\ n;\ell \vdash c_2}
            {n;\ell \vdash \texttt{if } e \texttt{ then } c_1
                           \texttt{ else } c_2}
\end{ruledisplay}
A \emph{command} is typed when $0;\mathsf{live} \vdash c$. Only
\textsc{Gate}, \textsc{Lift} and \textsc{Drop} name the temporary, and all
three need it live; so in a typed command nothing names the temporary after
its drop. The typed program of the body's affine calculus is the lift-free
command $\texttt{gate}\ U_1; \ldots; \texttt{gate}\ U_k; \drop\ t; \texttt{halt}$
that its elaboration lists.

\paragraph{The machine.} Fix a measurement record $o$, an infinite sequence over
$\{\mathsf{in}, \mathsf{out}\}$; the machine reads it from the left. A
configuration $\langle c; \psi; \eta; \sigma\rangle$ carries the command, the
joint state $\psi \in \R^{d} \otimes \R^{m}$, a store $\eta$ from parameters
to booleans, and the unread suffix $\sigma$ of $o$. With $P_{\mathsf{in}} =
P_S$, $P_{\mathsf{out}} = I - P_S$, $\beta(\mathsf{in}) = 1$ and
$\beta(\mathsf{out}) = 0$, the reductions are
\[
\begin{aligned}
 \langle \texttt{gate}\ U\,;c;\psi;\eta;\sigma\rangle
 &\longrightarrow \langle c;U\psi;\eta;\sigma\rangle,\\
 \langle \texttt{vgate}\ V\,;c;\psi;\eta;\sigma\rangle
 &\longrightarrow \langle c;(I \otimes V)\psi;\eta;\sigma\rangle,\\
 \langle \drop\ t;c;\psi;\eta;\sigma\rangle
 &\longrightarrow \langle c;\psi;\eta;\sigma\rangle,\\
 \langle \texttt{let }p_n=\lift(\meas\ t\in S)\texttt{ in }c;
          \psi;\eta;b\cdot\sigma\rangle
 &\longrightarrow
 \langle c;(P_b\otimes I)\psi;\eta[p_n\mapsto\beta(b)];\sigma\rangle,\\
 \langle \texttt{if } e \texttt{ then } c_1 \texttt{ else } c_2;
          \psi;\eta;\sigma\rangle
 &\longrightarrow \langle c_1;\psi;\eta;\sigma\rangle
   \ \text{if } \eta \models e, \ \text{else }
   \langle c_2;\psi;\eta;\sigma\rangle,
\end{aligned}
\]
and a configuration whose command is $\texttt{halt}$ is terminal. A
configuration is \emph{well typed} when $n;\ell \vdash c$ and $\eta$ binds
exactly $p_0, \ldots, p_{n-1}$. The initial configuration at visible input
$v$ is $\langle c; e_{i_0} \otimes v; \emptyset; o\rangle$.

\paragraph{The selected denotation.} By recursion on $c$, with a store $\eta$
and a record $\sigma$:
\[
\begin{aligned}
 E_{\sigma,\eta}(\texttt{halt}) &= I, \qquad
 E_{\sigma,\eta}(\texttt{gate}\ U\,;c) = E_{\sigma,\eta}(c)\,U, \\
 E_{\sigma,\eta}(\texttt{vgate}\ V\,;c) &= E_{\sigma,\eta}(c)\,(I \otimes V), \qquad
 E_{\sigma,\eta}(\drop\ t\,;c) = E_{\sigma,\eta}(c), \\
 \mathrlap{E_{b\cdot\sigma,\eta}(\texttt{let }p_n=\lift(\meas\ t\in S)\texttt{ in }c)} \\
   &= E_{\sigma,\eta[p_n\mapsto\beta(b)]}(c)\,(P_b \otimes I),\\
 E_{\sigma,\eta}(\texttt{if } e \texttt{ then } c_1 \texttt{ else } c_2)
   &= E_{\sigma,\eta}(c_1) \ \text{if } \eta \models e, \ \text{else }
     E_{\sigma,\eta}(c_2).
\end{aligned}
\]
The selected denotation of a typed command along $o$ is $E_o(c)(v) =
E_{o,\emptyset}(c)(e_{i_0} \otimes v)$. The drop of $c$ is
\emph{operationally safe} under $o$ when the map from visible inputs to
terminal joint states factors as $g \otimes Wv$.

\stmtMeta

\begin{proof}
Two parts, each a standard pair.

\paragraph{(1) Progress and preservation.}
\emph{Progress.} Let $\langle c;\psi;\eta;\sigma\rangle$ be well typed with
$c \ne \texttt{halt}$.
\begin{itemize}
\item a gate, a visible gate or a drop steps unconditionally;
\item a lift steps because $\sigma$ is a suffix of the infinite record $o$
  and so has a head;
\item a conditional steps because $\eta \models e$ is defined: \textsc{If}
  confines the parameters of $e$ to $p_0, \ldots, p_{n-1}$, which $\eta$
  binds.
\end{itemize}

\emph{Preservation.} The gate, visible-gate, drop and conditional steps keep
$n$ and $\eta$ and continue with a premise of the rule that typed $c$, at the
same $n$; the lift step continues with a command typed at $n+1$ and extends
$\eta$ by exactly $p_n$. So the store binds exactly the parameters in scope
after every step. After $\drop\ t$ the continuation is typed with the flag
$\mathsf{dropped}$, and no rule with that flag in its conclusion names the
temporary; so no later step names it.

\paragraph{(2) Termination and adequacy.}
\emph{Termination.} Every step replaces the command by a proper subcommand,
and commands are finite trees.

\emph{Adequacy.} By induction on the length of the run: if $\langle
c;\psi;\eta;\sigma\rangle$ runs to $\langle \texttt{halt};
\psi';\eta';\sigma'\rangle$ then $\psi' = E_{\sigma,\eta}(c)\,\psi$. For a
terminal configuration $E_{\sigma,\eta}(\texttt{halt}) = I$. Otherwise the
first step is one of the five reductions, each matching one clause of the
recursion:
\begin{itemize}
\item a gate contributes $U$, and a visible gate $I \otimes V$;
\item a drop contributes nothing;
\item a lift on outcome $b$ contributes the projection $P_b \otimes I$,
  together with the extended store;
\item a conditional selects the branch that $\eta$ selects.
\end{itemize}
The induction hypothesis applied to the continuation finishes. Starting from
$\langle c; e_{i_0} \otimes v; \emptyset; o\rangle$, the terminal state is
$E_{o,\emptyset}(c)(e_{i_0} \otimes v) = E_o(c)(v)$.

\paragraph{The last sentence.} It is a definition read twice: operational
drop safety asks whether $v \mapsto E_o(c)(v)$ factors, which is drop safety
of $E_o(c)$.
\end{proof}

\paragraph{The proportionality lemma.} The body's channel theorem and its
contextual-equivalence theorem both reduce to one linear-algebra fact.

\stmtProp
\begin{proof}
Write $k_{ib}$ and $m_b$ for the $b$-th rows. The $(b,b)$ entry of the
hypothesis reads $\sum_i \langle k_{ib}, x\rangle^{2} = \langle m_b,
x\rangle^{2}$ for all $x$; taking $x \perp m_b$ forces every $\langle k_{ib},
x\rangle = 0$, so $k_{ib} \in (m_b^{\perp})^{\perp} = \mathrm{span}(m_b)$:
$k_{ib} = c_{ib} m_b$, with $\sum_i c_{ib}^{2} = 1$ when $m_b \ne 0$. For two
nonzero rows $m_b, m_{b'}$ the $(b,b')$ entry gives $\sum_i c_{ib} c_{ib'}
\langle m_b, x\rangle \langle m_{b'}, x\rangle = \langle m_b, x\rangle
\langle m_{b'}, x\rangle$; the product of two nonzero linear forms is not the
zero polynomial, so $\sum_i c_{ib} c_{ib'} = 1$, and Cauchy--Schwarz between
the unit vectors $(c_{ib})_i$ and $(c_{ib'})_i$ holds with equality only when
they coincide. So $c_{ib} = c_i$ does not depend on $b$ over the nonzero rows
of $M$, and a zero row of $M$ is a zero row of every $K_i$; hence $K_i = c_i
M$.
\end{proof}

\paragraph{Contextual indistinguishability.} The body's theorem that the
criterion is contextual indistinguishability.

\stmtCtxEquiv
\begin{proof}
($\Rightarrow$) A weakenable program has the denotation of a separable one,
so no continuation distinguishes them. ($\Leftarrow$) Let $\sep(R,W)$ be
undistinguished, with $g = Re_{i_0}$ a unit vector, so its Kraus operators are
$g_i W$ with $\sum_i g_i^{2} = 1$. For a unit $u \ne e_0$, the reflection
exchanging $e_0$ and $u$ is in $F$ and satisfies $C^{\top}
e_0 = u$, so agreement of the two distributions at $b = 0$ says $\sum_i
\langle u, K_i x\rangle^{2} = \langle u, Wx\rangle^{2}$ for all such $u$ and
all $x$, and both sides are continuous in $u$, so it holds at $e_0$ as well. These are the quadratic forms of the symmetric matrices $\sum_i
(K_i x)(K_i x)^{\top}$ and $(Wx)(Wx)^{\top}$, which therefore coincide, and
\cref{lem:prop} gives $K_i = c_i W$. So the program is drop-safe, hence
weakenable by the body's criterion theorem.
\end{proof}

\paragraph{The summed channel.} The body's remark that an analysis reading
the channel a program implements cannot decide the drop rests on one pair of
commands.

\stmtChannel
\begin{proof}
Both apply $H$ to the temporary, reaching $\tfrac1{\sqrt2}(|0\rangle +
|1\rangle) \otimes v$. The first then applies $\mathrm{CZ}$, lifts on $\{0\}$,
and drops: $\mathrm{CZ}$ gives $\tfrac1{\sqrt2}(|0\rangle \otimes v + |1\rangle
\otimes Zv)$, so the branch selecting $0$ has denotation $\tfrac1{\sqrt2}
|0\rangle \otimes I$ and the other $\tfrac1{\sqrt2} |1\rangle \otimes Z$, one
Kraus operator each: both drops are safe, with branch maps $X \mapsto X/2$
and $X \mapsto ZXZ/2$. The second lifts on $\{0\}$, applies $t := t \oplus v$
(the visible bit controlling), and drops: the branch selecting $b$ keeps
$\tfrac1{\sqrt2}|b\rangle \otimes v$, and the controlled write sends it to
$\tfrac1{\sqrt2} \sum_c v_c\, |b \oplus c\rangle \otimes |c\rangle$, whose
Kraus operators are $\tfrac1{\sqrt2} P_0$ and $\tfrac1{\sqrt2} P_1$ in some
order, not proportional: both drops are unsafe, with branch map $X \mapsto
\operatorname{diag}(X)/2$ twice. Both pairs sum to the dephasing channel $X
\mapsto \operatorname{diag}(X)$, and every branch weight is
$\operatorname{tr}(X)/2$.
\end{proof}

\paragraph{A product at every input.} The body's remark that ``a product
at every input'' admits programs drop safety excludes is the following
dichotomy.

\stmtProductwise
\begin{proof}
For each $v$ with $Kv \ne 0$ write $Kv = g_v \otimes w_v$ with $g_v$ and
$w_v$ nonzero. Either all the $g_v$ are parallel or two of them are not, and
the two cases are settled separately.

\paragraph{All the temporary factors on one line.} If every $g_v$ lies on the
line of one vector $g$, put $W = (\langle g| \otimes I)\,K/\|g\|^{2}$. For
$Kv = g_v \otimes w_v = \lambda_v\,g \otimes w_v$ this gives $Wv = \lambda_v
w_v$ and $Kv = g \otimes Wv$; for $Kv = 0$ both sides vanish. So $K = g
\otimes W$.

\paragraph{Otherwise the visible factors are.} Choose $v_1, v_2$ with
$g_{v_1}, g_{v_2}$ linearly independent. By linearity
\[
  K(v_1 + v_2) \;=\; g_{v_1} \otimes w_{v_1} + g_{v_2} \otimes w_{v_2} ,
\]
which as a $d \times m$ matrix is $g_{v_1}w_{v_1}^{\top} +
g_{v_2}w_{v_2}^{\top}$, of rank two unless $w_{v_1}$ and $w_{v_2}$ are
parallel. A product vector, or zero, has rank at most one, so they are
parallel to one $w$. Any $v_3$ with $Kv_3 \ne 0$ has $g_{v_3}$ independent of
$g_{v_1}$ or of $g_{v_2}$, since otherwise $g_{v_1}$ and $g_{v_2}$ would be
parallel, and the same argument gives $w_{v_3}$ parallel to $w$. So every
$w_v$ lies on the line of $w$, and $G = (I \otimes \langle
w|)\,K/\|w\|^{2}$ is linear with $Kv = Gv \otimes w$ for every $v$, as
before.

\paragraph{The last sentence.} $K_i = (\langle i| \otimes I)K$ sends $v$ to
$(Gv)_i\,w$, so $K_i = w\,e_i^{\top}G$, and these are pairwise proportional
exactly when the rows of $G$ are, that is, when $G$ has rank at most one. If
$K$ is an isometry then $\|Gv\|\,\|w\| = \|v\|$ for every $v$, so $G$ is
injective and has rank $m$.
\end{proof}

\paragraph{Without typing.} The body's remark that the norm clause is the type
system's whole contribution to the criterion.

\stmtUntyped
\begin{proof}
The criterion theorem's proof uses the isometry hypothesis at one step, where
it makes $W$ orthogonal, and the norm clause supplies exactly that step. It
cannot be dropped, because a witness is a typed program and so denotes an
isometry: at $d = m = 1$ every map has a single Kraus operator and is
therefore drop-safe, while $v \mapsto 2v$ is not an isometry and no typed
program has its denotation.
\end{proof}

\paragraph{What the Gram matrix is.} Written over $\C$; over $\R$ read
$\dagger$ as $\top$ and $\bar U$ as $U$, and drop the second half of (1).

\stmtGram
\begin{proof}
All three parts are statements about one family: the Kraus operators $K_1,
\ldots, K_d$ of $K$, and the matrix $\Gamma$ they determine.

\paragraph{(1) Safety is the scalar condition.} Two directions, then the
reading over $\C$.

\emph{Safe $\Rightarrow$ scalar.} If $K = g \otimes W$ with $\|g\| = 1$ then
$W$ is an isometry of the visible register, hence unitary, and
$K_i^{\dagger}K_j = \bar g_i g_j I$.

\emph{Scalar $\Rightarrow$ safe.} Let $K_i^{\dagger}K_j = c_{ij}I$. Taking
traces, $m\,c_{ij} = \Gamma_{ij}$, so $c$ is Hermitian and positive
semidefinite; write $c = V D V^{\dagger}$ with $V$ unitary and $D$ diagonal,
and put $U = V^{\top}$. Rotating the temporary by $U$ diagonalises the
family:
\[
  K'_k \;=\; \sum_i U_{ki} K_i , \qquad
  K'^{\dagger}_k K'_l \;=\; \sum_{ij} \bar U_{ki} U_{lj} c_{ij} I
    \;=\; (\bar U c U^{\top})_{kl} I
    \;=\; (V^{\dagger} c V)_{kl} I \;=\; D_{kl} I .
\]
At most one diagonal entry is nonzero: if $\lambda_k$ and $\lambda_l$ were
both nonzero then $K'_k/\sqrt{\lambda_k}$ and $K'_l/\sqrt{\lambda_l}$ would
be isometries of $\C^{m}$, hence unitaries, with product $K'^{\dagger}_k
K'_l = 0$, which is impossible. So $K'_k = 0$ for all $k$ but one, $k_0$,
and
\[
  K \;=\; (U^{\dagger} \otimes I)K' \;=\; (U^{\dagger} e_{k_0}) \otimes K'_{k_0}
\]
is drop-safe.

\emph{The state reading, and why it is complex.} The temporary's final state
at a unit input $v$ has entries
\[
  \rho_t(v)_{ij} \;=\; \langle K_j v, K_i v\rangle
    \;=\; v^{\dagger} K_j^{\dagger} K_i v .
\]
Over $\C$ a matrix whose quadratic form $v \mapsto v^{\dagger}Av$ is constant
on unit vectors is a scalar, so $\rho_t$ is input-independent exactly when
every $K_j^{\dagger}K_i$ is a scalar, which is the first half. Over $\R$ this
fails. With $J$ the rotation of $\R^{2}$ by a right angle, take
\[
  K \;=\; \tfrac1{\sqrt2}\bigl(|0\rangle \otimes I + |1\rangle \otimes J\bigr).
\]
Since $v^{\top}Jv = 0$, it leaves the temporary in $\tfrac12 I$ at every real
unit input, while $I$ and $J$ are not proportional; at the complex input
$\tfrac1{\sqrt2}(1, i)$ the temporary is pure.

\paragraph{(2) The rank is the least temporary dimension.} Let $X$ be the $d
\times m^{2}$ matrix whose rows are the $K_i$ written out. Then $XX^{\dagger}
= \Gamma^{\top}$, so
\[
  \operatorname{rank}\Gamma \;=\; \operatorname{rank} X
    \;=\; \dim\operatorname{span}\{K_i\} ,
\]
which is the rank of the channel's Choi matrix and, as the two bounds below
show, the least number of Kraus operators in any representation of it---its
\emph{Choi rank}~\cite[p.~78]{Watrous2018TQI}.

\begin{itemize}
\item \emph{No smaller dimension works.} A typed program with a temporary of
  dimension $d'$ gives its channel $d'$ Kraus operators, so $d' \ge
  \operatorname{rank}\Gamma$.

\item \emph{That dimension is achieved.} Let $r = \operatorname{rank}\Gamma$
  and let $K'_1, \ldots, K'_r$ be Kraus operators of the channel, which exist
  because a channel always has a Kraus family of exactly that
  many~\cite{Choi1975CP,Watrous2018TQI}. The channel is trace preserving because $K$ is an
  isometry, so
  \[
    V v \;=\; \sum_k |k\rangle \otimes K'_k v
  \]
  is an isometry $\C^{m} \to \C^{r} \otimes \C^{m}$. Its columns extend to a
  unitary $U$ on $\C^{r} \otimes \C^{m}$ with $U(e_0 \otimes v) = Vv$, and
  the program that allocates at $e_0$, applies $U$ and drops has denotation
  $V$, temporary dimension $r$, and the channel of $K$.
\end{itemize}

\paragraph{(3) How the Gram matrix transports.} A visible gate $C$ after the
scope replaces $K_i$ by $CK_i$, and $\operatorname{tr}(K_i^{\dagger}
C^{\dagger} C K_j) = \Gamma_{ij}$. A gate $U$ on the temporary replaces $K_i$
by $\sum_j U_{ij} K_j$, and the Gram matrix of that family is $\bar U \Gamma
U^{\top}$, as computed in (1).
\end{proof}

The Gram matrix of a family of vectors is the standard invariant of the
family in quantum information: it decides linear independence, hence
unambiguous discrimination~\cite{Chefles1998Unambiguous}, and it is the
data of the family up to a unitary. Part (1) is the Knill--Laflamme
condition~\cite{Knill1997Theory} for the code that is the whole visible
register, with the $K_i$ as the errors and the temporary as the
environment; its rank-one coefficient matrix says the correction is
trivial, and the information-disturbance tradeoff of Kretschmann,
Schlingemann and Werner~\cite{Kretschmann2008Tradeoff} is its quantitative
form, a channel being close to correctable exactly when its complementary
channel is close to constant. The second half of (1) is noninterference in
the sense of Goguen and Meseguer~\cite{Goguen1982Security}, with the
visible input as the high input and the temporary as the low observer; Ying,
Feng and Yu define a quantitative version for quantum automata by the
distance between the observer's states along purged and unpurged
runs~\cite{Ying2013Noninterference}, and the impurity $\iota$ of
\cref{app:cost} is that degree for the drop.

\subsection{Clifford\texorpdfstring{$+T$}{+T} splits, in every dimension}\label{app:clifford}

\paragraph{The relativised criterion.} The body states the criterion over
a gate ring and defers its proof here.

\stmtRelativised
\begin{proof}
First the object both halves are about. If $U$ is drop-safe with unit final
state $g$, then $U_{(i,b),(0,c)} = g_i W_{bc}$, and row $0$ of the unitary
$W$ is a unit vector, so $\mathrm{proj}(U) = gg^{\dagger}$.

\paragraph{A witness makes the projector split.} A witness $\sep(R,W)$ over
$\mathcal A$ denotes $g \otimes W$ with $g = Re_0$, a column of a unitary
over $\mathcal A$. So $U$ is drop-safe and $\mathrm{proj}(U) = gg^{\dagger}$
splits, with $z = g$.

\paragraph{A splitting vector builds a witness.} Let $gg^{\dagger} =
zz^{\dagger}$ with $z$ over $\mathcal A$. Then $z$ is a unit vector on the
line of $g$, say $z = \lambda g$ with $|\lambda| = 1$, and the witness is
assembled in two pieces.

\begin{itemize}
\item \emph{The temporary factor.} By completability $z$ is a column of a
  unitary over $\mathcal A$; composing with a permutation matrix, which is
  over $\mathcal A$, gives a unitary $R$ over $\mathcal A$ with $Re_0 = z$.
  At $d = 2$ take $R = \bigl(\begin{smallmatrix} z_0 & -\bar z_1 \\ z_1 &
  \bar z_0 \end{smallmatrix}\bigr)$, over $\mathcal A$ because $\mathcal A$
  is closed under conjugation: every gate ring is completable at $d = 2$.

\item \emph{The visible factor.} Put
  \[
    W'_{bc} \;=\; \sum_i \bar z_i U_{(i,b),(0,c)}
      \;=\; \langle z, g\rangle W_{bc} \;=\; \bar\lambda W_{bc} ,
  \]
  which is $\bar\lambda W$ and so unitary, and which lies over $\mathcal A$
  because it is a polynomial in entries of $U$ and of $\bar z$.
\end{itemize}

Then $\sep(R,W')$ is typed, over $\mathcal A$, and denotes $z \otimes
\bar\lambda W = g \otimes W$, the denotation of $U$.
\end{proof}

Write $\mathcal R = \Z[1/\sqrt2, i] = \Z[\omega][1/\sqrt2]$ with $\omega =
e^{i\pi/4}$, and $\mathcal R_0 = \Z[1/\sqrt2]$ for its real elements. Three
facts about these rings are used. First, $\mathcal R$ is a principal ideal
domain: $\Z[\omega]$ is the ring of integers of $\Q(\zeta_8)$, whose class
number is one (a classical fact; the body gives the reference), and a
localisation of a principal ideal domain is one. Second, the units of
$\mathcal R_0$ are $\pm\sqrt2^{\,a}(1+\sqrt2)^{b}$ with $a, b \in \Z$: the
units of $\Z[\sqrt2]$ are $\pm(1+\sqrt2)^{b}$, and inverting the prime
$\sqrt2$ adds the one generator. Third, three unit norms: $|\sqrt2|^{2} = 2$,
$|1+\sqrt2|^{2} = (1+\sqrt2)^{2}$, and $|1+\omega|^{2} = 2 + \omega +
\bar\omega = 2+\sqrt2 = \sqrt2(1+\sqrt2)$, where $1+\omega$ is a unit of
$\mathcal R$ because $(2+\sqrt2)(2-\sqrt2) = 2$ and $2 = (\sqrt2)^{2}$ is a
unit. Write $\sigma$ for the
automorphism of $\Q(\omega)$ with $\sigma(\sqrt2) = -\sqrt2$; the Galois
group is abelian, so $\sigma$ commutes with complex conjugation and
$\sigma(|x|^{2}) = |\sigma(x)|^{2} \ge 0$.

\stmtCliffordT

\begin{proof}
Part (2) is part (1) fed to the relativised criterion, so the work is in (1).

\paragraph{(1) Every rank-one projection over $\mathcal R$ splits.} Let $P =
gg^{\dagger}$ with $\|g\| = 1$ have entries in $\mathcal R$. The splitting
vector is built in four steps.

\begin{enumerate}
\item \emph{A column, made primitive.} Some column $r = Pe_j = \bar g_j\,g$
  is nonzero; it lies in $\mathcal R^{d}$, and $P = rr^{\dagger}/\|r\|^{2}$.
  Divide $r$ by a greatest common divisor of its entries, which exists
  because $\mathcal R$ is a principal ideal domain: $r = \delta r'$ with $r'$
  primitive, and still $P = r'r'^{\dagger}/N$ for
  \[
    N \;=\; \|r'\|^{2} \;=\; \textstyle\sum_i |r'_i|^{2} \;\in\; \mathcal R_0 .
  \]

\item \emph{The norm is a unit.} $N P_{ik} = r'_i \bar r'_k$, so $N$ divides
  every product $r'_i \bar r'_k$; these products generate the ideal $(r'_0,
  \ldots, r'_{d-1})(\bar r'_0, \ldots, \bar r'_{d-1}) = (1)$. So $N$ is a
  unit of $\mathcal R$ and, being real, of $\mathcal R_0$: $N =
  \sqrt2^{\,a}(1+\sqrt2)^{b}$, positive.

\item \emph{The exponents have equal parity.} The conjugate $\sigma(N) =
  \sum_i |\sigma(r'_i)|^{2}$ is positive too, and $\sigma(N) =
  (-\sqrt2)^{a}(1-\sqrt2)^{b}$ has the sign of $(-1)^{a+b}$; so $a \equiv b
  \pmod 2$.

\item \emph{The unit whose norm is $N$.} Put $c = a \bmod 2$ and
  \[
    \varepsilon \;=\; \sqrt2^{\,(a-c)/2}\,(1+\sqrt2)^{(b-c)/2}\,(1+\omega)^{c},
  \]
  a unit of $\mathcal R$ with $|\varepsilon|^{2} = \sqrt2^{\,a-c}
  (1+\sqrt2)^{b-c}\bigl(\sqrt2(1+\sqrt2)\bigr)^{c} = N$.
\end{enumerate}

Then $z = r'/\varepsilon$ has entries in $\mathcal R$, $\|z\|^{2} =
N/|\varepsilon|^{2} = 1$, and $zz^{\dagger} = r'r'^{\dagger}/N = P$. At $d = 1$
the four steps degenerate rather than fail: the only rank-one projection on
$\C^{1}$ is $P = (1)$, so $r' = 1$, $N = 1$, $a = b = c = 0$ and $\varepsilon =
1$, giving $z = 1$.

\paragraph{(2) So the criterion is unchanged over $\mathcal R$.} A drop-safe
program over $\mathcal R$ has $\mathrm{proj}(U) = gg^{\dagger}$ over
$\mathcal R$, which splits by (1), and $\mathcal R$ is completable in every
dimension by the column lemma of Giles and Selinger cited in the body; the
body's relativised criterion then gives a witness over $\mathcal R$. The
converse is the criterion's other direction.
\end{proof}

\stmtRings

\begin{proof}
The proof of (1) used only three things about $\mathcal R$: that it is a
principal ideal domain closed under conjugation, that its real elements form
$\mathcal R_0$, and that every totally positive unit of $\mathcal R_0$ is
$|\varepsilon|^{2}$ for a unit $\varepsilon$ of $\mathcal R$. So the two
positive cases are a matter of checking those three.

\paragraph{The two rings that split.}
\begin{itemize}
\item \emph{$\Z[1/2, i] = \Z[i][1/2]$.} $\Z[i]$ has class number one, the
  real elements form $\Z[1/2]$, whose positive units are the $2^{a}$, and
  $2^{a} = |(1+i)^{a}|^{2}$ with $1+i$ a unit since $2 = -i(1+i)^{2}$.
\item \emph{$\Z[i/\sqrt2] = \Z[\sqrt{-2}][1/2]$.} $\Z[\sqrt{-2}]$ has class
  number one, the real elements again form $\Z[1/2]$, and $2^{a} =
  |\sqrt{-2}^{\,a}|^{2}$ with $\sqrt{-2}$ a unit since $2 = -(\sqrt{-2})^{2}$.
\end{itemize}
The three class numbers are one because the Minkowski bound of each of
$\Q(i)$ and $\Q(\sqrt{-2})$ is below $2$, and that of $\Q(\zeta_8)$ is below
$3$ while its only prime of norm $2$ is $(1-\omega)$. So the argument of (1)
goes through unchanged.

\paragraph{The two that do not.} Over a real ring $\mathcal A$ a splitting
vector $z$ of a real projection $P = gg^{\top}$ is $\lambda g$ with $\lambda
\in \mathcal A$ and $|\lambda| = 1$, so $z = \pm g$. That is what fails:
\begin{itemize}
\item over $\Z[1/2]$, the projection $P = \tfrac12\bigl(\begin{smallmatrix}
  1 & 1 \\ 1 & 1\end{smallmatrix}\bigr)$ has entries in the ring and $\pm(1,
  1)/\sqrt2$ does not;
\item over $\Z[1/\sqrt2]$, the projection $\tfrac12(I + (X+Z)/\sqrt2)$ has
  entries in the ring, and no unit vector on its ray does, by the computation
  below.
\end{itemize}
\end{proof}

\paragraph{The magic state, and the program that realises the separation.}
Take $g$, the $H$-type magic state $\cos\tfrac{\pi}{8}|0\rangle +
\sin\tfrac{\pi}{8}|1\rangle$, whose ray is $[(1+\sqrt2,\,1)]$ and whose
projector $\tfrac12(I + (X+Z)/\sqrt2)$ has entries in $\Z[1/\sqrt2]$. A unit
vector on that ray is $\lambda(1+\sqrt2, 1)$ with $\lambda^{2} =
(2-\sqrt2)/4$, whose norm down to $\Q$ is $1/8$, not a rational square, so no
such $\lambda$ lies in $\Q(\sqrt2)$, let alone in $\Z[1/\sqrt2]$, and the
projector does not split. Over $\mathcal R$ it does, with $\omega =
e^{i\pi/4}$: $z = \bigl(\tfrac{1+\omega}{2}, \tfrac{i(1-\omega)}{2}\bigr)$ has
$zz^{\dagger} = \tfrac12\bigl(I + \tfrac{X+Z}{\sqrt2}\bigr)$. One program
realises the separation: the single gate $U = \mathit{rot} \otimes
\mathit{rot}$, where $\mathit{rot}$ is the rotation of a qubit by $\pi/8$. Its
entries are $\cos^{2}\tfrac\pi8 = \tfrac{2+\sqrt2}{4}$, $\cos\tfrac\pi8
\sin\tfrac\pi8 = \tfrac{\sqrt2}{4}$ and $\sin^{2}\tfrac\pi8 =
\tfrac{2-\sqrt2}{4}$, all in $\Z[1/\sqrt2]$, though $\mathit{rot}$ alone is
over neither ring, and $U(|0\rangle \otimes v) = g \otimes \mathit{rot}\,v$,
so the drop is safe with the magic state as its final state. Over $\mathcal R$
this program is weakenable, contracting against $\bar z$ for the visible
factor $e^{-i\pi/8}\mathit{rot}$; over $\Z[1/\sqrt2]$ it is not.

\subsection{A command whose verdict depends on the outcome}\label{app:outcome}

\begin{planned}{the outcome-dependent command, and the fragment classification}
The proofs for the body's boundary section, and then the smallest temporary
that still obstructs, the fragments that escape, a family forcing a verdict to
read many record entries, and the closed form for consecutive lifts. Being
written; nothing here is claimed.
\end{planned}

\subsection{Several temporaries}\label{app:many}

\begin{planned}{composition, lifetimes, summaries, and borrows}
The proofs for the body's section on several temporaries, and then what a sound
procedure summary may forget, and how many test states certify a borrowed
temporary. Being written; nothing here is claimed.
\end{planned}

\subsection{The price of the verdict}\label{app:cost}

Throughout, a \emph{circuit program} is a program of the calculus whose
temporary and visible registers are qubits, $d = 2^{k}$ and $m = 2^{n}$, the
temporary allocated at $|0^{k}\rangle$, whose gates act on at most three
qubits of the joint register and are given by their matrices, each entry an algebraic number presented exactly by its minimal polynomial and an isolating rectangle in binary; its size counts the gates, the qubits and that presentation, and scalars are complex. $K$ is the denotation, an isometry $\C^{m} \to \C^{d}
\otimes \C^{m}$ with Kraus operators $K_i = (\langle i| \otimes I)K$. For a
unit vector $\psi$ of the visible register tensored with a reference register
$R$ of any finite dimension, $\rho_t(\psi) = \operatorname{Tr}_{vR}\bigl[(K
\otimes I_R)\,\psi\psi^{\dagger}\,(K \otimes I_R)^{\dagger}\bigr]$ is the state
of the temporary after the run; its \emph{impurity} $1 -
\operatorname{tr}\rho_t(\psi)^{2}$ vanishes exactly when it is pure, $\iota(K)
= \sup_{\psi}\bigl(1 - \operatorname{tr}\rho_t(\psi)^{2}\bigr)$ is its worst
case, and $\bar\iota(K)$ is its value at the maximally entangled input $\psi
= m^{-1/2}\sum_c |c\rangle \otimes |c\rangle$, whose visible marginal is the
maximally mixed state. The reference is essential: the swap $Kv = v \otimes
e_0$ of the body's productwise proposition leaves a pure temporary at every
pure input, yet at the maximally entangled input its temporary is maximally
mixed and $\iota(K) = \bar\iota(K) = 1 - 1/m$. Since $\rho_t(\psi) = \operatorname{Tr}_v[K
\sigma K^{\dagger}]$ depends on $\psi$ only through its visible marginal
$\sigma$, a reference of dimension $m$ suffices, the supremum is attained,
and $\bar\iota(K) \le \iota(K)$. Promise problems are meant throughout:
membership is in the promise version of the class, and hardness is under
polynomial-time many-one reductions of promise problems.

\paragraph{The temporary's state at the maximally entangled input.} It is
$\Gamma^{\top}/m$; hence $\bar\iota(K) = 1 -
\operatorname{tr}(\Gamma^{2})/m^{2}$, and $K$ is drop-safe if and only if
$\bar\iota(K) = 0$, if and only if $\iota(K) = 0$. Both propositions below
rest on this.

The joint state at the maximally entangled input is $(K \otimes I_R)\psi$,
and tracing out $R$ leaves $K K^{\dagger}/m$ on temporary and visible;
tracing out the visible register then gives $\operatorname{Tr}_v[KK^{\dagger}]
/ m$. Its $(i,j)$ entry is $\sum_{a,c} K_{(i,a),c}\,\overline{K_{(j,a),c}} =
\sum_{a,c} (K_i)_{ac} \overline{(K_j)_{ac}} = \operatorname{tr}(K_j^{\dagger}
K_i) = \Gamma_{ji}$, so the state is $\Gamma^{\top}/m$; its trace is
$\sum_i \operatorname{tr}(K_i^{\dagger}K_i) = \operatorname{tr}(K^{\dagger}K)
= m$ divided by $m$, as it must be, and its purity is
$\operatorname{tr}(\Gamma^{2})/m^{2}$. For the equivalences: a state of
trace one is pure exactly when it is a rank-one projector, so $\bar\iota(K)
= 0$ says $\Gamma$ has rank one, which is form (4) of the body's definition
of drop safety. If $K = g \otimes W$ with unit $g$
then for every input, with or without a reference, the joint state is $g
\otimes (W \otimes I_R)\psi$ and $\rho_t(\psi) = gg^{\dagger}$, so $\iota(K)
= 0$; and $\iota(K) = 0$ implies $\bar\iota(K) = 0$, which closes the cycle.
(Since $\Gamma$ is positive semidefinite with trace $m$, rank one is also
the identity $\Gamma^{2} = m\Gamma$; the checker takes the rank, and the
suite checks the two against each other.)


\stmtClassicalCost
\begin{proof}
Every gate in $\{X, \mathrm{CNOT}, \mathrm{Toffoli}\}$ permutes the joint
computational basis, so a program over them does too. Write its denotation as
\[
  K e_x \;=\; |a(x)\rangle \otimes |b(x)\rangle ,
\]
where $x$ ranges over visible basis inputs, $a(x)$ is the basis state the
temporary is left in and $b(x)$ the one the visible register is left in. By
the body's definition the drop is safe exactly when the temporary factor is
one vector for every $x$, that is, exactly when
\emph{$a$ is constant}. Both halves below are about that one function.

\paragraph{Membership.} A pair $x \ne y$ with $a(x) \ne a(y)$ is a
certificate of unsafety, and it is polynomially checkable. These programs
contain no measurement, so a program is one list of gates applied in order,
with no branching and nothing to select. A joint basis state is a bit
string --- the visible positions holding $x$, the allocated ones their
initial values --- and each of the three gates reads at most three positions
of that string and rewrites one: $X$ flips its own bit, $\mathrm{CNOT}$ flips
its target when the control bit is $1$, and $\mathrm{Toffoli}$ flips its
target when both control bits are $1$. Each gate is therefore a
constant-time update of the string, so for a program of $s$ gates $a(x)$ ---
the temporary's positions once all $s$ have been applied --- costs $O(s)$.
Safety is the absence of such a pair, so the problem is in coNP.

\begin{figure}
\centering
\includegraphics{figures/classical-price}
\caption{Look at the mirror. Every gate drawn permutes the computational
basis, so the temporary is left in one basis state $a(x)$ for each input, and
every gate of the compute block has a partner in the uncompute block, so the
scratch qubits reach the drop at $|0\rangle$ on every input. The copy into the
target is the one gate without a partner, and the only mark on the target's
wire; that is why the temporary reaches the drop holding $B(x)$ and nothing
else, and why the drop is safe exactly when $B$ is constant.}
\label{fig:classical-price}
\end{figure}

\paragraph{Hardness.} We reduce from unsatisfiability. The drop sees
constancy of $a$, and what must be decided is unsatisfiability, so the
reduction takes three steps: evaluate $B$ reversibly, read $a$ off the
resulting program, and convert constancy into unsatisfiability with one
evaluation. \Cref{fig:classical-price} draws the reduction on a two-gate $B$.

Let $B$ be a Boolean circuit over $\{\wedge, \neg\}$ with $n \ge 1$ inputs and
$s$ gates; a CNF formula converts to such a circuit in polynomial size. Taking
$n \ge 1$ is free: a formula with no variables is decided by one evaluation,
and the visible register below is the one carrying $B$'s inputs, so it must be
nonempty for the fixed unsafe program of step 3 to exist.

\begin{enumerate}
\item \emph{Evaluate $B$ reversibly}, by Bennett's
  compute--copy--uncompute~\cite{Bennett1973Logical}. Allocate $s$ scratch
  qubits and one target, all at $|0\rangle$, and in topological order replace
  each gate of $B$ by a reversible one writing into a \emph{fresh} scratch
  qubit:
  \begin{itemize}
  \item $\wedge$ with inputs on wires $u, w$ becomes a Toffoli from $u, w$
    into the fresh scratch qubit;
  \item $\neg$ with input on wire $u$ becomes a CNOT from $u$ into the fresh
    scratch qubit, followed by $X$ on it.
  \end{itemize}
  Then a CNOT copies the circuit's output wire into the target, and then the
  whole gate sequence is applied again in reverse order. Every gate used is an
  involution, so the reverse pass returns each scratch qubit to $|0\rangle$ on
  every input $x$; only the target keeps a value. The net map is
  \[
    |x\rangle \otimes |0^{s}, 0\rangle \;\longmapsto\;
    |x\rangle \otimes |0^{s}, B(x)\rangle .
  \]

\item \emph{Read the verdict off it.} Take the temporary to be the register of
  $s+1$ qubits (the scratch and the target), allocated at $|0^{s+1}\rangle$,
  and the visible register to be $x$. This is a circuit program of size
  $O(n+s)$, and by step 1 its denotation is
  \[
    K e_x \;=\; |0^{s}, B(x)\rangle \otimes e_x ,
    \qquad\text{so}\qquad a(x) = (0^{s},\,B(x)) .
  \]
  So $a$ is constant exactly when $B$ is, and the drop is safe if and only if
  $B$ is constant.

\item \emph{Convert constancy into unsatisfiability.} A constant $B$ may be
  constantly $1$, so constancy is not yet unsatisfiability. One evaluation
  closes the gap. Compute $B(0^{n})$ in polynomial time and output:
  \begin{itemize}
  \item if $B(0^{n}) = 1$: a fixed unsafe program---a CNOT from the first
    visible qubit into a temporary qubit. Here $B$ is satisfiable, and the
    output is unsafe.
  \item if $B(0^{n}) = 0$: the program of step 2. It is safe exactly when $B$
    is constant, and $B(0^{n}) = 0$ makes that constant $0$; so it is safe
    exactly when $B \equiv 0$, that is, when $B$ is unsatisfiable.
  \end{itemize}
\end{enumerate}

In both cases the program produced is drop-safe if and only if $B$ is
unsatisfiable, and the reduction runs in polynomial time. So the problem is
coNP-hard, and with the first half, coNP-complete.
\end{proof}

With the $s+1$ temporaries kept as separate names in the calculus of
the body's section on several temporaries, the same thing holds per name: each
scratch qubit is separately constant and the target is separately constant
exactly when $B$ is constant. The predicate is not the one of Liu, Zhou and
Meng~\cite{Liu2026Uncomputation}, who ask whether some unitary can restore the
ancilla while preserving the data and prove that coNP-complete over
$\{\mathrm{Toffoli}, X\}$ by a reduction from satisfiability through
Jordan's strong equivalence of reversible
circuits~\cite{Jordan2014StrongEquivalence}. The two predicates separate on
the program above: it stores $B(x)$ in the target, so it is uncomputable in
their sense for every $B$, by running backwards, and drop-safe in ours only
for constant $B$.

The exact verdict, with no promise on the size of the leak, inherits
\cref{thm:classicalcost} over any gate set containing $X$, CNOT and Toffoli,
and lies in PSPACE for gates over a fixed number field, since each entry of
$K$ is a sum of exponentially many products of gate entries, computable in
polynomial space, and form (3) of the body's definition enumerates the minors;
Tanaka shows the exact form of non-identity check to be
NQP-complete~\cite{Tanaka2009ExactNonIdentity}, and where the exact drop
verdict sits between coNP and PSPACE is open.

\stmtAverageCost
\begin{proof}
The algorithm estimates one number, $\operatorname{tr}\rho^{2}$ for $\rho =
\Gamma^{\top}/m$, and reads the promise off it.

\paragraph{The algorithm.} A uniformly random basis input $e_x$ leaves the
temporary in $\operatorname{Tr}_v[K e_x e_x^{\dagger} K^{\dagger}]$, whose
average over $x$ is
\[
  \operatorname{Tr}_v[KK^{\dagger}]/m \;=\; \Gamma^{\top}/m \;=\; \rho ,
\]
the temporary's state at the maximally entangled input; two runs on
independently drawn inputs leave the two temporaries in $\rho \otimes \rho$. The swap test on two registers in state $\rho \otimes \sigma$
accepts with probability $\tfrac12(1 +
\operatorname{tr}(\rho\sigma))$~\cite{Buhrman2001Fingerprinting}, here
\[
  \tfrac12(1 + \operatorname{tr}\rho^{2}) \;=\; 1 - \tfrac12\bar\iota(K) .
\]
Repeat $O(1/\varepsilon^{2})$ times and accept when the empirical acceptance
rate exceeds $1 - \varepsilon/4$: by Hoeffding's inequality the two cases,
rate $1$ and rate at most $1 - \varepsilon/2$, are told apart with constant
error. Every step runs the program---whose gates the machine
applies directly, which is why the fixed finite gate set is
hypothesised---or a swap test on polynomially many qubits. For
gates presented as arbitrary algebraic matrices the bound needs a compilation
step with its own error analysis, which we do not carry out.

\paragraph{The reflection named in the body.} For $m \ge 2$, $U = I -
2|w\rangle\langle w|$ has nonzero $+1$ and $-1$ eigenspaces, and the two
readings of the leak come apart:
\begin{itemize}
\item an input superposing a $+1$ with a $-1$ eigenvector makes $c = \langle
  \psi|(U \otimes I)|\psi\rangle$ vanish, the temporary maximally mixed and
  $\iota(K_U) = \tfrac12$;
\item at the maximally mixed input $c = \operatorname{tr}(U)/m = 1 - 2/m$, so
  $\bar\iota(K_U) = \tfrac12(1 - (1 - 2/m)^{2}) = 2/m - 2/m^{2}$.
\end{itemize}
\end{proof}

\begin{planned}{the price on a restricted gate set}
The gate sets on which the verdict is cheap, and what a gate one step outside
such a set costs. Being written; the statement is not in the body either.
\end{planned}

% ===========================================================================
% B.  THE CATALOGUE
% ---------------------------------------------------------------------------
% What the body summarises in one paragraph: the seventeen programs, the two
% modelled producer rules, the criterion, and the runtime-check table for
% the witness. Reference material, and allowed to read like it.
% ===========================================================================
\section{The catalogue}\label{app:catalogue}

\subsection{Weakenable without being separable: three witnesses}
\label{app:witness}

The body's \emph{weakening witness} asks that \emph{some} typed program of
the shape $\sep(R,W)$ have the same denotation, not that the program be of
that shape. The
difference is easy to state and easy to lose, so here are three catalogue
programs beside the witness the body's criterion produces for each. The middle
column is the denotation the two share.

\begin{center}
\small
\begin{tabular}{@{}L{.30\linewidth}L{.30\linewidth}L{.30\linewidth}@{}}
\toprule
\textbf{the program} & \textbf{its denotation} & \textbf{a witness $\sep(R,W)$}\\
\midrule
\code{silq-two-gate}: \code{t := H t; t := H t; drop t}
  & $|0\rangle \otimes v$
  & $R = I$, $W = I$ --- allocate and drop, touching nothing\\[2pt]
\code{dj-constant}: the body's Deutsch--Jozsa text, with its joint $U_f$
  & $|{-}\rangle \otimes HD_fHv$
  & $R = HX$, carrying $|0\rangle$ to $|{-}\rangle$; $W = HD_fH$, which is
    $\pm I$\\[2pt]
\code{grover-oracle}: compute $f$ into the ancilla, kick the phase back,
  uncompute
  & $|0\rangle \otimes D_f v$
  & $R = I$; $W = D_f = \operatorname{diag}((-1)^{f(x)})$, a legal gate
    because it is real and orthogonal\\
\bottomrule
\end{tabular}
\end{center}

None of the three programs \emph{is} $\sep(R,W)$: each has a gate that touches
both registers, and the first has two gates on the temporary where the shape
allows one. All three are weakenable, and by \cref{thm:ctxequiv} no
continuation distinguishes any of them from the witness beside it. Two
catalogue programs, \code{in-shape} and \code{do-nothing}, are literally of the
shape; the other ten safe ones are in the position of these three.

Each witness is computed the way the criterion's proof computes it: read $g$
and $W$ off the denotation, then prepare $g$ from $e_{i_0}$ by the reflection
through $(e_{i_0}-g)^{\perp}$, which for $g = |{-}\rangle$ and $e_{i_0} =
|0\rangle$ is $HX$ up to sign.

The seventeen catalogue programs of the body, with the verdict of the
criterion and of the two modelled producer rules. \textbf{R1} accepts when
every gate is a permutation; \textbf{R2} adds that no gate writes visible data
into the temporary. Computed by \code{dqt catalogue} and pinned by three
tests: one asserting R1 unsound, one asserting R2 sound and strictly
incomplete, and one asserting the totals printed here.
The names are the checker's own and only \code{example-B} is one of the body's
running programs: it is program~(b). Program~(a) is \code{park} and
program~(c) is \code{tilt-read}; \code{example-A} and \code{example-C} are
two further programs, not those.

\begin{center}
\small
\begin{tabular}{@{}l c c c@{}}
\toprule
program & R1 & R2 & criterion \\
\midrule
\code{example-A}      & accept & accept & safe \\
\code{example-B}      & refuse & refuse & unsafe \\
\code{example-C}      & refuse & refuse & safe \\
\code{park}          & refuse & refuse & safe \\
\code{silq-two-gate} & refuse & refuse & safe \\
\code{tilt-read}     & refuse & refuse & unsafe \\
\code{in-shape}      & refuse & refuse & safe \\
\code{dj-constant}   & refuse & refuse & safe \\
\code{dj-balanced}   & refuse & refuse & safe \\
\code{bv-0}          & refuse & refuse & safe \\
\code{bv-1}          & refuse & refuse & safe \\
\code{grover-oracle} & accept & refuse & safe \\
\code{kickback}      & refuse & refuse & unsafe \\
\code{copy-copy}     & refuse & refuse & safe \\
\code{cz-copy}       & refuse & refuse & unsafe \\
\code{compute-alone} & \textbf{accept} & refuse & \textbf{unsafe} \\
\code{do-nothing}    & accept & accept & safe \\
\midrule
\textbf{accepts}     & 4 & \textbf{2} & \textbf{12} \\
\bottomrule
\end{tabular}
\end{center}

Of the twelve safe programs, the ray each leaves the temporary on, and the
verdict of the restore-to-allocation rule (\code{dqt pool}):

\begin{center}
\small
\begin{tabular}{@{}lll@{}}
\toprule
program & left on & restore rule \\
\midrule
\code{example-A}, \code{silq-two-gate}, \code{grover-oracle}, \code{do-nothing} & $[1,0]$ & accepts \\
\code{example-C}, \code{park}, \code{in-shape}, \code{copy-copy} & $[1,1]$ & \textbf{refuses} \\
\code{dj-constant}, \code{dj-balanced}, \code{bv-0}, \code{bv-1} & $[1,-1]$ & \textbf{refuses} \\
\bottomrule
\end{tabular}
\end{center}

The runtime-check table emitted by \code{dqt residual} for the witness command
of the body's boundary section: on the branch selecting $\{0,1\}$,
\textsc{refuse}; on the branch selecting $\{2\}$, \textsc{admit}. On the same command with the
read deleted, \textsc{admit} on both.


\bibliographystyle{ACM-Reference-Format}
\bibliography{refs}
\end{document}
